\section{Findings: How Lab and Field \rev{Scientific Practitioners} Currently Use AI Tools}
\label{sec:usecases}

In the following three sections, we first examine how \rev{scientific practitioners working in lab and field settings} currently incorporate AI tools into their work (\sec{sec:usecases}). We then describe the barriers they identified that limit adoption in practice (\sec{sec:barriers}). Finally, we present five archetypal speculative designs for future AI assistants that \rev{participants} came up with, which illustrate how they envision, in the ideal case, AI supporting both their day-to-day and long-term research (\sec{sec:spec-design}). Below we present two categories of current AI use cases: at the desk, and beyond the desk.

%\vspace{-0.4em}

\subsection{AI Tools for Knowledge Work `At the Desk'}

Of the 12 \rev{participants} we interviewed, all had heard of AI tools, particularly recent Generative AI (genAI) tools. However, AI usage patterns varied widely, ranging from little-to-no use (P4, P7, P9) to sporadic reliance (P1, P5, P6, P12) to using on a daily basis (P2, P3, P8, P10, P11). The most commonly-mentioned AI tool was ChatGPT, followed by other text-based LLMs such as Claude, Google Gemini, and Perplexity AI-powered search.

For most participants we interviewed, AI use remained peripheral---used intermittently on the computer (i.e., `at their desk') at the beginning or end stages of their workflows, such as during experimental planning, data analysis, or manuscript writing---and was contextually disconnected from their day-to-day physical work in the lab and the field.

The most prevalent AI use cases were for programming and debugging (P1, P2, P3, P5, P6, P8, P10, P12), followed by brainstorming and ideation around experiment protocols (P1, P2, P3, P8, P10, P11). Several participants used AI for clarifying and organizing their thoughts (P3, P10, P11), while others had integrated or considered integrating it into their scholarly communication practices through paper writing assistance (P2, P4, P7, P12), figure generation (P2, P10, P11), and visualization support for presentations (P10). Information-seeking behaviors included using AI for literature review tasks (P1, P8), personal tutoring for learning new concepts (P3, P5, P11), and general knowledge queries (P1, P2, P3, P6, P7, P8, P10, P11, P12). Other use cases included interpreting technical manuals (e.g. setting up trail cameras) (P8), verifying experimental calculations (P5, P12), writing emails, or assisting with paperwork (P7, P11, P12).

When participants walked the interviewer through their work-spaces, many drew clear distinctions between computational and physical spaces. P1, who works in a wet lab studying the neural activity of rodent behavior through \textit{in vivo} electrophysiology, stated that he \textit{``mainly only uses [AI] for data analysis.''} P1 only uses AI tools on his computer at home or at the computer station in an office next to the wet lab. P5, who works in the same field, also expressed that he \textit{``uses [AI] seldomly. I probably don't utilize it as much as I could for my benefit, but it's because it's mainly for coding.''}

Note that both of these \rev{participants} spend significant amounts of time directly interfacing with animals: training, feeding, socializing them, surgically implanting custom-built ``drives'' to measure neural behavior, and running experiments. In these contexts, where research spans years and involves hands-on, dynamic interaction with living systems in-situ, AI tool utility was perceived as limited and more suited to purely computational work back at their desks.

In sum, the majority of \rev{participants}' current AI usage was for computer-based knowledge work at their desks, which confirms findings from the three known prior studies of scientists and AI use~\cite{morris2023scientists,wagman2025generative,obrien2025how}, along with broader studies of other types of knowledge workers~\cite{guo2025from, khojah2024beyond, pu2025assistance, liu2023code}. Where our study goes beyond prior work is uncovering use cases, barriers, and design opportunities for scientific AI use `beyond the desk,' which we present below.


\subsection{AI Tools for Physical Work `Beyond the Desk'}


Several described using AI tools in their lab or field environments, such as voice-activated assistance when their hands were busy (P8), help during DIY (Do-It-Yourself) engineering tasks (P8), and hardware troubleshooting (P2, P3, P11).

For instance, P8, a field researcher who studies primates (e.g., macaques, orangutans, chimpanzees) in remote locations around the world such as rainforests, often needs to engineer experimental setups on-site to conduct his research. These apparatuses cannot be pre-constructed due to their size and are sometimes built with salvaged or improvised materials on-site. P8 recounted a logistically-challenging experience in the Republic of Congo, where he had to construct an experimental apparatus from \textit{``driftwood that [he] found and random objects that [he] uses.''} P8 once asked ChatGPT’s voice assistant mode to help him construct an apparatus made of acrylic, wood, and epoxy to observe primate inhibitory control: \textit{``I can directly ask [ChatGPT] questions, like trivial stuff when I build something, which I'm not great at. I sort of figure stuff out on the go, but if I ask it, `Oh, I need to cut acrylic glass, what tool should I use, what's the right blade?' ... for even stuff like that, I can use it in the field.''} This use case illustrates how AI tools can help \rev{participants} \textit{``figure out tools and solutions''} (P8) and serve as low-friction, just-in-time support for improvisational DIY problem-solving.

In another case, P11, a \rev{doctoral student researcher in the field of experimental psychology} works with virtual reality and game-controller equipment in her experiments and uses ChatGPT to troubleshoot hardware problems. She once nearly canceled a day of scheduled study participants due to issues with a computer graphics card \textit{``but then I gave it to ChatGPT, [and asked] what are all the possible issues for this? What are some of the workarounds, can I replace the cable? Like, going through all the different diagnostics I can do on my end before I take that equipment to a specialist or get outsourced help.''} She later said that she found this form of AI assistance empowering: \textit{``I'm working on a computer vision project right now, which I have zero expertise in, but with ChatGPT, I feel I'm more empowered to do stuff. I don't have to know every single thing.}''

These examples show how \rev{scientific practitioners} need to pick up a wide range of improvised DIY skills on-the-job, often in domains outside their formal academic training. In P8's case, although his expertise was in primate observation, he pointed out that he needed domain knowledge in mechanical engineering and woodwork. For P11, although she \rev{conducts research in the field of cognitive psychology}, in order to run her experiments, she needed to have domain knowledge in computer hardware.